% auto-generated by extract_acquisition.py -- do not hand-edit
\textbf{Parameter} & \textbf{LIDC} \\
\midrule
Manufacturer set & GE MEDICAL SYSTEMS (670), SIEMENS (201), Philips (73), TOSHIBA (66) \\
Tube voltage (kVp), typical & 120.0 \\
Full-dose effective mAs, median (range) & 230 (0--11200) \\
Pitch, typical & 0.94 \\
Slice thickness (mm), thin series & 0.60--5.00 \\
Reconstruction kernel(s) & STANDARD, BONE, B30f, LUNG \\
% --- caption (paste into manuscript.tex, replacing the old one) ---
\caption{Acquisition parameters. v0.5 contains LIDC only (n = 1{,}010); the LIDC cells are regenerated from \texttt{metadata.json} by the committed extraction script \texttt{analysis/extract\_acquisition.py} (build provenance: 1010 files, combined SHA-256 \texttt{20476a8105319e6192eebf233e274a4f6f5f787a574bd605809963f3160df870}). The AAPM 2016 and Mayo LDCT-PD columns are the v1.0 roadmap plan and are not part of v0.5. Full per-scan values --- including the exact scanner-model string, pitch, kernel, and per-scan effective mAs --- are exposed by the loader without imposing a single value. LIDC's broad multi-vendor coverage (four manufacturers: GE, Siemens, Philips, Toshiba) is a noteworthy strength for cross-vendor generalization studies.}
